\section{Protocol Overview and Threat Model}
\label{sec:protocol}

\subsection{Payment Channel Lifecycle}
\label{sec:payment-channel}

In a payment channel, two parties broadcast an on-chain funding transaction.
Once the channel is open, they transact off-chain by exchanging signed
commitments, and update the state by exchanging new ones. When the state is
updated, the parties exchange revocation secrets for the previous commitment
and revoke the old state. If one party publishes an old commitment after
revocation, the other may claim all funds. Between state updates, any
commitment that has not yet been revoked is a valid unilateral close: if either
party publishes one there is no penalty. Punishment is triggered only when a
revoked commitment is published.

Channels close either cooperatively, with both parties agreeing on the final
balance and settling in a single chain transaction, or unilaterally. In a
unilateral close one party broadcasts an unrevoked commitment, and both must
wait for a relative timelock (CSV) to expire. That delay is what gives the
counterparty time to detect whether the published state was revoked. If it was,
the counterparty punishes the cheater and claims the entirety of the channel
funds.

\subsection{Multi-Hop HTLC Routing}
\label{sec:multi-hop-htlc-routing}

When sender and recipient have no direct channel, the payment is routed through
intermediaries along a path $S \to F_1 \to \dots \to F_n \to R$. Because
intermediaries are untrusted, each hop is conditioned on a single shared
secret, so that the payment either completes along the entire path or completes
nowhere.

\medbfparagraph{Lock round}
The receiver generates a preimage $x$ and sends the sender an invoice
containing the hash $y = H(x)$ and the requested amount $v$, keeping $x$
secret. The sender adds an HTLC to its channel with $F_1$, locking funds
conditional on $y$ before a deadline $d_0$; after $d_0$ the payer may reclaim
them. The unlock condition is to present $x'$ with $H(x') = y$ before the
deadline. Each forwarder in turn locks an HTLC on its outgoing channel under
the \emph{same} hash, deducting a routing fee, so that amounts decrease and
deadlines tighten toward the receiver:
\[
\small
v_0 > v_1 > \cdots > v_n = v,
\qquad
d_0 > d_1 > \cdots > d_n .
\]

No funds move during this round. Each HTLC becomes an additional conditional
output in the two parties' current commitment transaction, and is therefore
subject to the revocation mechanism of \Cref{sec:payment-channel}. Locking is
safe for $F_i$ precisely because it holds an incoming HTLC on the same hash for
a strictly larger amount: if its outgoing HTLC is claimed, the secret needed to
claim its incoming one is thereby exposed.

\medbfparagraph{Release round}
The receiver reveals $x$ to the last forwarder to claim the final HTLC. Since
every HTLC shares the preimage, each intermediary claims on its incoming
channel using the same $x$, learning it when the next hop claims downstream. To
claim $v_{n-1}$ from $F_{n-1}$, $F_n$ presents $x$ on its incoming HTLC and
retains $v_{n-1} - v_n$ as its fee. The preimage propagates backward hop by hop
until the sender's HTLC is settled. If any party fails to forward and the
timelock expires, funds return to whoever locked them. If honest forwarders
claim on time, everyone is paid; otherwise no one is.

\medbfparagraph{Why deadlines must decrease}
Consider $F_i$ holding an incoming HTLC at $d_{i-1}$ and an outgoing one at
$d_i$. If $F_{i+1}$ claims the outgoing HTLC at the latest permitted moment
$d_i$, then $F_i$ learns $x$ at $d_i$ and retains the window $(d_i, d_{i-1})$
in which to settle its incoming claim. The gap $d_{i-1} - d_i$ is the CLTV
delta, and it is precisely the time reserved for $F_i$ to recover its funds.
Were the ordering inverted, so that $d_{i-1} < d_i$, then $F_i$'s incoming HTLC
would expire and be refunded before the downstream party claims at $d_i$: $F_i$
would have paid $v_i$ downstream while recovering nothing upstream, losing the
full forwarded amount. \Cref{sec:minimality} derives the positivity of this
delta and the non-emptiness of the resulting window from a concrete block
clock, rather than assuming them.

\medbfparagraph{Residual liveness requirement}
A non-empty claim window is necessary but not sufficient: nothing in the
timelock structure compels $F_i$ to use it. An intermediary offline during
$(d_i, d_{i-1})$ forfeits its incoming HTLC even though every party followed
the protocol and no one cheated. Routing security therefore reduces, at this
point, to a liveness assumption on honest intermediaries, which corresponds in
deployment to watchtower infrastructure~\cite{mccorry2019pisa}.

Furthermore, the entire release round depends on the receiver disclosing $x$.
While funds are locked, every party on the path is waiting for one of two
events:

\begin{itemize}
  \item the receiver releases $x$, which propagates backward so that each party
  claims its HTLC and collects its fee; or
  \item the deadline $d_i$ passes, the HTLC times out, and the party recovers
  its funds.
\end{itemize}

No one can force the receiver to reveal $x$. A receiver that simply declines
immobilises the capital of every party on the path until the deadlines elapse
--- a griefing denial of service. Separately, the \emph{wormhole attack} occurs
when two dishonest, non-adjacent parties communicate directly, relaying $x$
between themselves and skipping the intermediaries in between without
disclosing it to them. The skipped intermediary is honest throughout: it
forwards value and earns no fee. Both pathologies are reachable in our model.

\subsection{Threat Model and Security Goals}
\label{sec:threat}

We assume the standard Dolev--Yao adversary model~\cite{dolev1983security}. The
adversary controls the whole network: it can read, modify, and inject messages.
Its only limitation is that it cannot break cryptographic primitives. We
additionally model long-term key compromise as a separate adversarial
capability, so that properties can be stated conditionally on honesty rather
than assuming it globally.

\medbfparagraph{Liveness assumptions}
Security at the routing layer degrades into a liveness requirement: a node that
goes offline loses funds even in an otherwise correct execution. We encode this
with four timing restrictions, distinguishing passive orderings from active
obligations.

\begin{itemize}
  \item \textbf{Structural clock constraints.} T1 (\code{RedeemBeforeTimeout}),
  T2a (\code{ForwardTimeGap}), and T2b
  (\code{HonestPartiesActBeforeIncomingTimeout}) are assumed structural
  constraints of the CLTV clock.
  \item \textbf{Active claiming.} T3 (\code{IntermediaryMustClaim}) requires an
  intermediary to act: it must claim its incoming HTLC once its outgoing HTLC
  is redeemed.
\end{itemize}

T1, T2a and T2b carry timing soundness; T3 carries intermediary safety.
\Cref{sec:minimality} shows that none is implied by the others. Our security
goals translate directly into verified lemmas: payment atomicity, intermediary
safety, routing causality and authentication, timing safety, and value
conservation.
